\documentclass[11pt]{article}
\usepackage[margin=1in]{geometry}
\usepackage[T1]{fontenc}
\usepackage{lmodern}
\usepackage{amsmath}
\usepackage{amssymb}
\usepackage[colorlinks=true,linkcolor=blue,citecolor=blue,urlcolor=blue]{hyperref}
\usepackage[nameinlink,noabbrev]{cleveref}

\title{Advanced Display Mathematics Stress Paper}
\author{Stage Two Corpus}
\date{}

\begin{document}
\maketitle

\begin{abstract}
This fixture concentrates the display-mathematics constructs that most often
break a LaTeX-to-OMML conversion path. It exercises aligned multi-line
equations with embedded text, manually tagged equations, piecewise
definitions, several matrix environments, labelled arrows, simultaneous
overbraces and underbraces, binomial coefficients, fractions, and indexed
sums. The surrounding prose is written so that a text-difference metric over
the document body remains meaningful, while the mathematical surface remains
the primary object under test.
\end{abstract}

\section{Estimator and Objective}
\label{sec:objective}

Consider a regression problem in which the parameter vector
$\theta \in \mathbb{R}^{p}$ is estimated from $n$ paired observations
$(x_i, y_i)$ with $x_i \in \mathbb{R}^{p}$ and $y_i \in \mathbb{R}$. The
penalised objective combines a quadratic loss with a ridge term, and we write
the gradient on a second line so that the alignment marker is exercised
directly:
\begin{align}
L(\theta)
  &= \sum_{i=1}^{n} \bigl(y_i - x_i^{\top}\theta\bigr)^2
     + \lambda \lVert \theta \rVert_2^2,
     && \text{with } \lambda \geq 0, \label{eq:loss}\\
\nabla L(\theta)
  &= -2 X^{\top}(y - X\theta) + 2\lambda\theta. \label{eq:grad}
\end{align}
The first-order condition obtained by setting \cref{eq:grad} to zero yields the
closed-form minimiser, which we display with a manual tag rather than an
automatic number so that the tag-handling branch of the converter is reached:
\begin{equation}
\stepcounter{equation}% keep the hyperref anchor distinct from the auto-numbered displays
\hat{\theta}(\lambda)
  = \bigl(X^{\top}X + \lambda I_p\bigr)^{-1} X^{\top} y. \tag{$\star$}
\end{equation}
The expression in \eqref{eq:loss} reduces to ordinary least squares when
$\lambda = 0$, in which case the estimator inherits the familiar projection
geometry of the column space of $X$.

\section{A Piecewise Decision Rule}
\label{sec:cases}

The fitted parameter induces a classifier through a thresholded score. We write
the decision rule as a piecewise expression, which forces the converter to map
a \texttt{cases} environment onto the corresponding stacked structure in the
target format:
\begin{equation}
d(x) =
\begin{cases}
1, & \text{if } x^{\top}\hat{\theta} \geq \tau,\\[2pt]
0, & \text{if } 0 \leq x^{\top}\hat{\theta} < \tau,\\[2pt]
-1, & \text{otherwise},
\end{cases}
\label{eq:cases}
\end{equation}
where $\tau$ is a fixed cutoff chosen on held-out data. The middle branch of
\cref{eq:cases} contains a compound inequality, and the textual annotations
inside each branch test the preservation of upright text within an otherwise
mathematical display.

\section{Covariance and Design Matrices}
\label{sec:matrices}

Two matrix environments appear below. The first is a symmetric correlation
matrix written with square brackets, and the second is a small design block
written with round brackets so that both delimiter families are exercised
within the same section:
\begin{equation}
\Sigma =
\begin{bmatrix}
1 & \rho & \rho^2\\
\rho & 1 & \rho\\
\rho^2 & \rho & 1
\end{bmatrix},
\qquad
B =
\begin{pmatrix}
1 & x_{1} & x_{1}^2\\
1 & x_{2} & x_{2}^2\\
1 & x_{3} & x_{3}^2
\end{pmatrix}.
\label{eq:matpair}
\end{equation}
The determinant of the leading principal block is reported through a
determinant environment, which carries vertical-bar delimiters and a multi-row,
multi-column body:
\begin{equation}
\begin{vmatrix}
1 & \rho\\
\rho & 1
\end{vmatrix}
= 1 - \rho^2 .
\label{eq:det}
\end{equation}
The matrices in \cref{eq:matpair} and the determinant in \cref{eq:det} share the
same scalar $\rho$, so a reader can verify the entries against one another.

\section{Annotated Transformations}
\label{sec:annotated}

Conversion paths frequently mishandle decorations that sit above or below a
sub-expression. The following display combines a labelled arrow with a
simultaneous overbrace and underbrace, so that both vertical decorations and the
arrow annotation are evaluated together:
\begin{equation}
\underbrace{X^{\top}X}_{\text{Gram matrix}}
\;\xrightarrow{\;+\,\lambda I_p\;}\;
\overbrace{X^{\top}X + \lambda I_p}^{\text{regularised}} .
\label{eq:arrow}
\end{equation}
The transformation in \cref{eq:arrow} is the operation that guarantees
invertibility for every $\lambda > 0$, because the added diagonal shifts every
eigenvalue upward by exactly $\lambda$.

\section{Combinatorial and Summation Notation}
\label{sec:sums}

A final display gathers a binomial coefficient, a nested fraction, and an
indexed sum into one numbered line so that the converter must place all three
operators within a single equation body:
\begin{equation}
\binom{n}{k}
= \frac{n!}{k!\,(n-k)!},
\qquad
S_n = \sum_{i=1}^{n} \frac{1}{i\,(i+1)}
    = 1 - \frac{1}{\,n+1\,}.
\label{eq:binom}
\end{equation}
The telescoping identity in \cref{eq:binom} provides a closed form for the
partial sum, and the binomial coefficient on the left supplies an independent
combinatorial expression in the same display.

\section{An Inline-Heavy Paragraph}
\label{sec:inline}

The remaining prose is deliberately dense with inline mathematics so that the
converter must interleave short formulas with ordinary words at high frequency.
For a centred predictor the leverage of observation $i$ equals
$h_{ii} = x_i^{\top}(X^{\top}X)^{-1}x_i$, which satisfies
$0 \leq h_{ii} \leq 1$ and $\sum_{i=1}^{n} h_{ii} = p$. When $\lambda$ grows the
effective degrees of freedom
$\mathrm{df}(\lambda) = \operatorname{tr}\!\bigl[X(X^{\top}X+\lambda I_p)^{-1}X^{\top}\bigr]$
decrease monotonically from $p$ toward $0$, so the map
$\lambda \mapsto \mathrm{df}(\lambda)$ is strictly decreasing on
$(0, \infty)$. A unit increase in the $j$th coordinate of $x$ shifts the score
by $\hat{\theta}_j$, and the standardised statistic
$z_j = \hat{\theta}_j / \widehat{\operatorname{se}}(\hat{\theta}_j)$ is compared
against the quantile $\Phi^{-1}(1-\alpha/2)$ at level $\alpha$. These inline
fragments, taken together with the numbered displays in
\cref{sec:objective,sec:cases,sec:matrices,sec:annotated,sec:sums}, give a
compact but demanding test of mathematical fidelity.

\end{document}
